\documentclass{aastex701}

\usepackage{amsmath,amssymb}

\newcommand{\dd}{\mathrm{d}}
\newcommand{\Mdot}{\dot{M}}
\newcommand{\Menv}{M_{\rm env}}
\newcommand{\Mcore}{M_{\rm core}}
\newcommand{\Mp}{M_{\rm p}}
\newcommand{\Rrcb}{R_{\rm rcb}}
\newcommand{\Rp}{R_{\rm p}}
\newcommand{\Rs}{R_{\rm s}}
\newcommand{\Tint}{T_{\rm int}}
\newcommand{\Teq}{T_{\rm eq}}
\newcommand{\cs}{c_{\rm s}}
\newcommand{\Lint}{L_{\rm int}}
\newcommand{\Lbol}{L_{\rm bol}}
\newcommand{\Lxuv}{L_{\rm XUV}}
\newcommand{\Ladv}{L_{\rm adv}}
\newcommand{\Fxuv}{F_{\rm XUV}}
\newcommand{\orchard}{\textsc{orchard}}

\shorttitle{ORCHARD Mass Loss Methods}
\shortauthors{Tejada Arevalo}

\begin{document}

\title{Atmospheric Mass Loss in the ORCHARD Planetary Evolution Code}

\author{Roberto Tejada Arevalo}
\email{rarevalo@ucsc.edu}
\affiliation{Department of Astrophysical Sciences, Princeton University, Princeton, NJ 08544, USA}

\begin{abstract}
We describe the atmospheric mass-loss framework implemented in the
\orchard{} planetary evolution code.  The module couples three physical
mechanisms---isothermal Parker winds, bolometric and intrinsic
energy-limited escape, and XUV photoevaporation---within a two-phase
evolutionary picture consisting of an early boil-off stage and a
subsequent post-boil-off stage.  Mass loss is self-consistently coupled
to the interior structure.  Envelope mass is removed by uniformly scaling
the Lagrangian mass grid.  The hydrostatic equilibrium solver then
recomputes the density, pressure, radius, and temperature profiles.  We
detail each mass-loss rate prescription, the phase-transition criterion,
the advective energy feedback, grid-adjustment procedure, and
timestep constraints.
\end{abstract}

\keywords{planets and satellites: atmospheres --- planets and satellites:
physical evolution --- methods: numerical}

% ======================================================================
\section{Introduction}\label{sec:intro}
% ======================================================================

Close-in sub-Neptune and super-Earth exoplanets lose their
hydrogen-helium envelopes through a combination of core-powered mass
loss \citep{Ginzburg2016, Gupta2019} and XUV photoevaporation
\citep{Lopez2013, Owen2017}.  This sculpts the observed radius valley
separating bare rocky cores from planets that retain volatile envelopes.

\orchard{} implements a unified mass-loss module that captures both
the early, rapid boil-off phase---in which the planet sheds mass on a
timescale shorter than its Kelvin--Helmholtz contraction
time---and the subsequent post-boil-off phase dominated by XUV-driven
escape.  The treatment follows \citet{Tang2024}.  It is coupled at every
evolutionary timestep to the code's Henyey relaxation solver for
hydrostatic equilibrium and to its implicit transport solver for entropy
and composition.

Section~\ref{sec:formulation} presents the shrinking-domain Lagrangian
formulation that justifies adding mass loss as a clean operator-split
step on top of the existing structure equations.
Section~\ref{sec:rates} presents the individual mass-loss rate
prescriptions.  Section~\ref{sec:phases} describes the two-phase logic
and the transition criterion.  Section~\ref{sec:grid} explains how
the Lagrangian grid is adjusted after mass removal and how
hydrostatic equilibrium is restored.
Section~\ref{sec:diagnostics} summarizes the diagnostic outputs.

% ======================================================================
\section{Shrinking-Domain Lagrangian Formulation}\label{sec:formulation}
% ======================================================================

The structure equations of \citet{TejadaArevalo2026} (their
Equations~1--6) describe a planet of fixed total mass $\Mp$ on the
Lagrangian domain $M_{r}\in[0,\Mp]$.  Mass loss makes the outer boundary
time-dependent: the planet's enclosed mass becomes
\begin{equation}\label{eq:Mt}
  M(t) \;=\; \Mp(t_{0})
  \;-\; \int_{t_{0}}^{t}\Mdot_{\rm wind}(t')\,\dd t',
  \qquad \Mdot_{\rm wind}(t)\ge 0,
\end{equation}
and the structure equations live on the shrinking Lagrangian domain
$M_{r}\in[0,M(t)]$.  We choose the absolute Lagrangian coordinate
$M_{r}$ as the independent variable rather than a normalised coordinate
$\xi = M_{r}/M(t)$.  At fixed interior $M_{r}$ no advective term enters
the time derivative.  A normalised coordinate, in contrast, would inject
$(M_{r}\Mdot/M^{2})\,\partial/\partial M_{r}$ corrections into hydrostatic
equilibrium, energy, and composition equations.  This is the
formulation adopted by \texttt{MESA} \citep{Paxton2011, Paxton2019} for
stellar winds and by every \texttt{MESA}-based planet
photoevaporation paper \citep[e.g.,][]{Lopez2013, Owen2017,
ChenRogers2016, Mordasini2020, Malsky2020, Rogers2021}.

The consequences for Equations~1--6 of \citet{TejadaArevalo2026} are:

\begin{enumerate}
  \item \textbf{Hydrostatic equilibrium and mass continuity (their
        Eqs.~1--2) are unchanged in form} on $[0,M(t)]$:
        \begin{equation}\label{eq:hse}
          \frac{\dd P}{\dd M_{r}}
            = -\frac{G\,M_{r}}{4\pi\,r^{4}}
              + \frac{\Omega(t)^{2}}{6\pi\,r}\,,
          \qquad
          \frac{\dd r}{\dd M_{r}}
            = \frac{1}{4\pi\,r^{2}\,\rho}\,,
        \end{equation}
        because at fixed $M_{r}$ both expressions are quasi-static
        statements of force and mass balance, independent of
        $\Mdot$ to leading order.  In particular, the Henyey
        relaxation solver of \orchard{} requires no new equations
        or new state variables; its 2N block-tridiagonal Jacobian
        is unaffected.

  \item \textbf{The outer pressure boundary condition is re-evaluated
        at the current $(M(t),R(t))$.}  Surface gravity becomes
        $g_{\rm surf}(t) = G\,M(t)/R(t)^{2}$.  The atmospheric
        boundary condition module returns $T_{\rm int}(t)$ and
        $T_{\rm eff}(t)$ at the post-mass-loss surface conditions.
        \orchard{}'s atmosphere modules (\texttt{c23}, \texttt{c26},
        \texttt{f07}, \texttt{rock\_gray}, etc.) are queried unchanged.

  \item \textbf{The interior energy and composition equations (their
        Eqs.~3--6) are unchanged in form for surviving shells.}  At
        fixed Lagrangian $M_{r}$, the entropy equation
        \begin{equation}\label{eq:Lprime}
          \frac{\partial L}{\partial M_{r}}
            = -\,T\,\frac{\partial S}{\partial t}\bigg|_{M_{r}}
              -\sum_{i}\!\bigl(\partial U/\partial X_{i}\bigr)_{S,\rho}\,
                \frac{\partial X_{i}}{\partial t}\bigg|_{M_{r}}
              + \epsilon_{\rm rad} + \epsilon_{L}\,,
        \end{equation}
        and the composition equations
        \begin{equation}\label{eq:Xprime}
          N_{A}\,\frac{\partial X_{i}}{\partial t}\bigg|_{M_{r}}
            = -\,\frac{\partial}{\partial M_{r}}
              \!\bigl(4\pi\,r^{2}\,F_{i}\bigr)\,,
        \end{equation}
        retain their original right-hand sides.  No $\Mdot$ term
        appears in the interior PDE; this is what allows the
        transport solver in \orchard{} to be re-used unchanged.

  \item \textbf{A new global equation (mass conservation) accompanies
        the moving boundary:}
        \begin{equation}\label{eq:dMdt}
          \frac{\dd M}{\dd t} \;=\; -\,\Mdot_{\rm wind}(t).
        \end{equation}

  \item \textbf{The surface composition is advected away with the
        wind.}  Defining the species mass budget
        $M_{i}(t) = \int_{0}^{M(t)}\!X_{i}\,\dd M_{r}$ and applying
        Leibniz' rule on the moving upper limit gives
        \begin{equation}\label{eq:Xboundary}
          \frac{\dd M_{i}}{\dd t}
            = -\!\left[4\pi\,r^{2}\,F_{i}\right]_{0}^{M(t)}
              \;-\; X_{i,{\rm surf}}\,\Mdot_{\rm wind}.
        \end{equation}
        Because the wind is launched from the upper atmosphere where
        convection has homogenised composition, $X_{i,{\rm surf}}$ is
        the cell-zero composition of the post-transport profile.
        No new flux term enters the transport solver; the surface
        composition is simply discarded with the lost mass during
        the rezone step (Section~\ref{sec:grid}).

  \item \textbf{The wind carries away thermal, kinetic, and
        gravitational energy.}  Defining the planet's total
        energy $E_{\rm tot}(t) = \int_{0}^{M(t)}\!
        \bigl[U + \tfrac{1}{2}v^{2} - GM_{r}/r\bigr]\,\dd M_{r}$ and
        applying Leibniz' rule on the upper limit gives
        \begin{equation}\label{eq:Eledger}
          \frac{\dd E_{\rm tot}}{\dd t}
            = \int_{0}^{M(t)}\!\!
              \bigl(\epsilon_{\rm rad}+\epsilon_{L}\bigr)\,\dd M_{r}
              - L_{\rm surf}
              - \Mdot_{\rm wind}\!
                \left[h_{\rm surf} + \tfrac{1}{2}v_{\rm wind}^{2}
                      - \frac{G\,M}{R}\right],
        \end{equation}
        where $h = U + P/\rho$ is specific enthalpy.  The bracketed
        term is the specific total energy of the escaping material.
        For energy-limited photoevaporation it is implicit in the
        $\Mdot$ formula, since the heating efficiency $\eta$ already
        encodes how much XUV energy goes into wind kinetic energy
        versus re-radiation.  It is nevertheless tracked explicitly in
        \orchard{} as a diagnostic (Section~\ref{sec:diagnostics}).
\end{enumerate}

In this formulation, mass loss is a clean operator-split addition to
the existing solver.  The transport and Henyey relaxation steps run
unchanged on the current domain.  Afterwards $\Mdot_{\rm wind}$ is
evaluated from one of the prescriptions (Section~\ref{sec:rates}); the
mass grid is updated to remove $\Delta M = \Mdot_{\rm wind}\,\Delta t$
of envelope mass (Section~\ref{sec:grid}); the boundary condition is
re-evaluated at the new $(M(t+\Delta t),R(t+\Delta t))$; and a final
hydrostatic re-equilibration absorbs the residual.

% ======================================================================
\section{Mass-Loss Rate Prescriptions}\label{sec:rates}
% ======================================================================

Three independent mass-loss rate estimates are computed at each timestep.
The final rate is determined by the evolutionary phase
(Section~\ref{sec:phases}).

% ----------------------------------------------------------------------
\subsection{Isothermal Parker Wind}\label{subsec:parker}
% ----------------------------------------------------------------------

Following \citet{Parker1958} and \citet{Tang2024} (their Equations~4--6),
we solve the isothermal, spherically symmetric, steady-state wind
equation for the velocity at the radiative--convective boundary (RCB).

The isothermal sound speed is
\begin{equation}\label{eq:cs}
  \cs = \sqrt{\frac{k_{\rm B}\,\Teq}{\mu}}\,,
\end{equation}
where $\Teq$ is the equilibrium temperature and $\mu$ is the mean
molecular weight of the outflowing gas (default
$\mu = 2.35\,m_{\rm H}$).  The sonic radius, at which the wind
velocity equals the sound speed, is
\begin{equation}\label{eq:Rs}
  \Rs = \frac{G\,\Mp}{2\,\cs^{2}}\,.
\end{equation}

The Parker wind equation relates the velocity $v$ at radius $r$ to the
sound speed through the implicit relation
\begin{equation}\label{eq:parker}
  \left(\frac{v}{\cs}\right)^{2}
  - \ln\left(\frac{v}{\cs}\right)^{2}
  = -3 + 4\ln\!\left(\frac{r}{\Rs}\right) + \frac{4\,\Rs}{r}\,.
\end{equation}
For subsonic flow at the RCB ($\Rrcb < \Rs$), the velocity is obtained
by inverting Equation~\eqref{eq:parker} using the Lambert-$W$ function
\citep[see][]{Cranmer2004}
\begin{equation}\label{eq:lambert}
  \left(\frac{v_{\rm rcb}}{\cs}\right)^{2}
  = -W_{0}\!\bigl(-e^{-f}\bigr),
  \qquad
  f = 3 - 4\ln\!\left(\frac{\Rrcb}{\Rs}\right) - \frac{4\,\Rs}{\Rrcb}\,,
\end{equation}
where $W_{0}$ is the principal branch of the Lambert-$W$ function.  A
physical solution requires the argument $-e^{-f}$ to lie in
$[-1/e,\,0)$.  If $\Rrcb \ge \Rs$, the flow is already supersonic
and $v_{\rm rcb} = \cs$ is adopted as a lower bound.

The Parker wind mass-loss rate is then
\begin{equation}\label{eq:mdot_parker}
  \Mdot_{\rm parker} = 4\pi\,\Rrcb^{2}\,\rho_{\rm rcb}\,v_{\rm rcb}\,,
\end{equation}
where $\rho_{\rm rcb}$ is the density at the RCB.  This rate is
\emph{not} energy-limited; it depends on the local thermodynamic
conditions at the RCB rather than on the total available luminosity.

% ----------------------------------------------------------------------
\subsection{Energy-Limited Escape}\label{subsec:elim}
% ----------------------------------------------------------------------

Two energy-limited rates cap the mass loss when the available luminosity
is insufficient to sustain the hydrodynamic wind.

\paragraph{Bolometric energy-limited rate.}
The total luminosity intercepted at the RCB---including both the
intrinsic luminosity $\Lint$ and re-radiated stellar heating---sets an
upper bound on the escape rate \citep[Equation~22 of][]{Tang2024}:
\begin{equation}\label{eq:mdot_bol}
  \Mdot_{\rm bol} = \frac{\Lbol^{\rm rcb}}{G\,\Mp / \Rrcb}\,,
\end{equation}
where the bolometric luminosity at the RCB is
\begin{equation}\label{eq:Lbol}
  \Lbol^{\rm rcb} = \Lint + 4\pi\,\Rrcb^{2}\,\sigma_{\rm SB}\,\Teq^{4}\,.
\end{equation}

\paragraph{Intrinsic energy-limited rate.}
When only the planet's own cooling luminosity drives the outflow
\citep[Equation~27 of][]{Tang2024}:
\begin{equation}\label{eq:mdot_int}
  \Mdot_{\rm int} = \frac{\Lint}{G\,\Mp / \Rrcb}\,.
\end{equation}
An optional bottleneck correction replaces the denominator with
$G\Mp/\Rrcb - G\Mp/R_{\tau_h}$, where $R_{\tau_h}$ is the radius at
optical depth $\tau = 2/3$.  This correction is not used in the default
configuration.

% ----------------------------------------------------------------------
\subsection{XUV Photoevaporation}\label{subsec:xuv}
% ----------------------------------------------------------------------

After the boil-off phase, XUV-driven escape dominates the mass-loss
budget for sub-Neptunes.  \orchard{} implements the photoevaporation
prescription of \citet{ChenRogers2016} (their \S2.3, hereafter CR16).
CR16 combine the energy-limited \emph{escape} rate of
\citet{Erkaev2007} with the recombination-limited \emph{radiative}
ceiling of \citet{MurrayClay2009}, and select the smaller of the two
at each timestep.

\paragraph{Stellar XUV luminosity.}  Two parametrisations of the host
star's XUV evolution are supported.  The default
\citep[\texttt{xuv\_evolution\_model = ribas05},][]{Ribas2005} is a
saturation plateau plus power-law decay,
\begin{equation}\label{eq:Lxuv}
  \Lxuv(t) =
  \begin{cases}
    f_{\rm sat}\,L_{\star}
      & t \le t_{\rm sat}\,,\\[4pt]
    f_{\rm sat}\,L_{\star}
      \left(\dfrac{t}{t_{\rm sat}}\right)^{-1.5}
      & t > t_{\rm sat}\,,
  \end{cases}
\end{equation}
with $f_{\rm sat} = 10^{-3.5}$ and $t_{\rm sat} = 100$~Myr.  The
alternative \citep[\texttt{xuv\_evolution\_model = rogers21},
following][]{Rogers2021} adds an explicit stellar-mass dependence,
\begin{equation}\label{eq:Lxuv_rogers}
  \frac{\Lxuv}{L_{\rm bol}}
    = 10^{-3.5}\!\left(\frac{M_{\star}}{M_{\odot}}\right)^{\!-0.5}
      \!\!\!\times
      \begin{cases}
        1 & t < t_{\rm sat}\\[4pt]
        \bigl(t/t_{\rm sat}\bigr)^{-1-a_{0}} & t \ge t_{\rm sat}
      \end{cases}
\end{equation}
with $a_{0} = 0.5$ and
$t_{\rm sat} = 100\,(M_{\star}/M_{\odot})^{-1}$~Myr.  The XUV flux at
the planet's orbital distance $a$ is then
\begin{equation}\label{eq:Fxuv}
  \Fxuv = \frac{\Lxuv}{4\pi\,a^{2}}\,.
\end{equation}

\paragraph{Energy-limited rate (CR16 Eq.~1).}  \citet{Lopez2013} adopt
a bare $\Rp^{3}$ geometric factor.  CR16 replace it with the more physical
$\Rp\,R_{\rm EUV}^{2}$ form---the EUV is \emph{absorbed} at the
photoionisation radius $R_{\rm EUV}$ but the binding energy per gram
is set at the optical photosphere $\Rp$---and include the tidal
correction of \citet{Erkaev2007},
\begin{equation}\label{eq:mdot_xuv_en}
  \Mdot_{\rm XUV}^{\rm en\text{-}lim}
    = \frac{\eta_{\rm XUV}(v_{\rm esc})\,
            \pi\,F_{\rm EUV}\,\Rp\,R_{\rm EUV}^{2}}
           {G\,\Mp\,K_{\rm tidal}(\xi)},
\end{equation}
where $\xi = R_{\rm Hill}/\Rp$ and the Erkaev tidal factor is
\begin{equation}\label{eq:Ktidal}
  K_{\rm tidal}(\xi) = 1 - \frac{3}{2\xi} + \frac{1}{2\xi^{3}}\,.
\end{equation}
For $\xi \to \infty$ (isolated planet) $K_{\rm tidal}\to 1$ and
Equation~\eqref{eq:mdot_xuv_en} reduces to the standard
energy-limited form; for close-in planets $K_{\rm tidal} < 1$ and
mass loss is enhanced.

\paragraph{Velocity-dependent heating efficiency.}  Following the
hydrodynamic simulations of \citet{OwenJackson2012} as parametrised by
\citet{Rogers2021},
\begin{equation}\label{eq:eta_xuv}
  \eta_{\rm XUV}(v_{\rm esc})
    = \eta_{0}\!\left(\frac{v_{\rm esc}}{15\,{\rm km\,s^{-1}}}\right)^{\!\!-\alpha_{\eta}},
  \qquad v_{\rm esc} = \sqrt{2\,G\,\Mp/\Rp}\,,
\end{equation}
with default $\eta_{0} = 0.1$ and $\alpha_{\eta} = 2$.  Setting
\texttt{eta\_xuv\_velocity\_dependent = False} recovers the legacy
constant-$\eta$ form.

\paragraph{Photoionisation radius.}  $R_{\rm EUV}$ is defined as the
radius at which the photoionisation optical depth reaches unity.
Following \citet{MurrayClay2009} and CR16 (their Eq.~2),
\begin{equation}\label{eq:Reuv}
  R_{\rm EUV} \approx \Rp + H\,
    \ln\!\left(\frac{P_{\rm photo}}{P_{\rm EUV}}\right),
  \qquad H = \frac{k_{B}\,T_{\rm photo}}{\mu_{\rm atm}\,g_{\rm surf}}\,,
\end{equation}
where the photoionisation base pressure is set by
$\sigma_{\nu_{0}}\,n_{\rm H}\,H = 1$, giving
\begin{equation}\label{eq:Peuv}
  P_{\rm EUV}
    = \frac{m_{\rm H}\,G\,\Mp}{\sigma_{\nu_{0}}\,\Rp^{2}},
  \qquad
  \sigma_{\nu_{0}} = 6\!\times\!10^{-18}
    \!\left(\frac{h\nu_{0}}{13.6\,{\rm eV}}\right)^{\!-3}\!{\rm cm^{2}},
\end{equation}
with default photon energy $h\nu_{0} = 20$~eV.  $R_{\rm EUV}$
typically exceeds $\Rp$ by 10--50\% for sub-Neptunes; the implied
geometric enhancement of $\sim 1.2\text{--}2.3\times$ in
Equation~\eqref{eq:mdot_xuv_en} is the leading correction over the
bare $\Rp^{3}$ form.

\paragraph{Recombination-limited rate (Murray--Clay 2009).}  At
$F_{\rm EUV}\gtrsim 10^{4}\,{\rm erg\,s^{-1}\,cm^{-2}}$, Lyman-$\alpha$
radiative cooling balances photoionisation heating, the wind
temperature saturates at $T_{\rm wind}\!\sim\!10^{4}$~K, and the
constant-$\eta$ assumption breaks.  The full
\citet{MurrayClay2009} expression (their Eq.~20; CR16 Eq.~3) is
\begin{equation}\label{eq:mdot_rr}
  \Mdot_{\rm XUV}^{\rm rr\text{-}lim}
    = \pi\!\left(\frac{G\,\Mp}{c_{s,{\rm rr}}^{2}}\right)^{\!2}
      c_{s,{\rm rr}}\,m_{\rm H}\,
      \!\left(\frac{F_{\rm EUV}\,G\,\Mp}
                   {h\nu_{0}\,\alpha_{\rm rec}\,
                    R_{\rm EUV}^{2}\,c_{s,{\rm rr}}^{2}}\right)^{\!\!1/2}
      \exp\!\left(2 - \frac{G\,\Mp}{c_{s,{\rm rr}}^{2}\,R_{\rm EUV}}\right),
\end{equation}
where $c_{s,{\rm rr}} = \sqrt{2\,k_{B}\,T_{\rm wind}/m_{\rm H}}$ is the
sound speed of the fully-ionised wind, the factor of 2 accounting for
protons and electrons.  This sound speed is not to be confused with the
bolometric $c_{s} = \sqrt{k_{B}\,\Teq/\mu}$ used by the boil-off Parker
wind.  Here $\alpha_{\rm rec} = 2.7\!\times\!10^{-13}\,{\rm cm^{3}\,s^{-1}}$
is the case-B recombination coefficient at $T_{\rm wind} = 10^{4}$~K.
The exponential factor encodes the gravitational suppression of the
density at the wind launch radius and is retained in full---it varies
by orders of magnitude across the parameter space and must not be
linearised.

\paragraph{Composite photoevaporation rate.}  The instantaneous
photoevaporation rate is the smaller of the two regimes (CR16 \S2.3
selection rule):
\begin{equation}\label{eq:mdot_photoev}
  \Mdot_{\rm photoev}
    = \min\!\bigl(\Mdot_{\rm XUV}^{\rm en\text{-}lim},\;
                   \Mdot_{\rm XUV}^{\rm rr\text{-}lim}\bigr).
\end{equation}
Both rates are evaluated and saved as diagnostics.

% ======================================================================
\section{Two-Phase Evolution}\label{sec:phases}
% ======================================================================

The mass-loss evolution is divided into two phases, following
\citet{Tang2024}.

% ----------------------------------------------------------------------
\subsection{Boil-Off Phase}\label{subsec:boiloff}
% ----------------------------------------------------------------------

During boil-off the planet cannot contract fast enough to radiate
away the energy deposited by the escaping wind.  The instantaneous
rate is set by the \emph{minimum} of the three mechanisms:
\begin{equation}\label{eq:mdot_boiloff}
  \Mdot_{\rm boiloff}
  = \min\!\bigl(\Mdot_{\rm parker},\;
                \Mdot_{\rm bol},\;
                \Mdot_{\rm int}\bigr).
\end{equation}

Two characteristic timescales govern the transition out of this phase:
\begin{enumerate}
  \item The \textbf{mass-loss timescale}
        (Equation~7 of \citealt{Tang2024}):
        \begin{equation}\label{eq:tM}
          t_{M} = \frac{\Menv}{\Mdot}\,.
        \end{equation}
  \item The \textbf{Kelvin--Helmholtz contraction timescale}
        (Equation~9 of \citealt{Tang2024}):
        \begin{equation}\label{eq:tcool}
          t_{\rm cool}
          = \frac{G\,\Mcore\,\Menv}{\alpha\,\Rrcb\,L_{\rm total}}\,,
        \end{equation}
        where $\alpha \le 1$ is an efficiency parameter (default
        $\alpha = 1$) and $L_{\rm total} = \Lint + L_{\rm irr}$ is the
        total luminosity including irradiation.
\end{enumerate}

\subsection{Phase Transition}\label{subsec:transition}

The planet exits the boil-off phase when the contraction timescale
drops below the mass-loss timescale:
\begin{equation}\label{eq:transition}
  t_{M} \ge t_{\rm cool}\,.
\end{equation}
Once this condition is met (and $t > 0$), the code transitions
irreversibly to the post-boil-off phase.

% ----------------------------------------------------------------------
\subsection{Core-Powered Mass Loss}\label{subsec:cpml}
% ----------------------------------------------------------------------

In the post-boil-off regime, an additional channel can drive escape
even when the stellar XUV flux is exhausted.  The planet's own
internal cooling luminosity launches a transonic Bondi-like wind from
above the RCB.  This \emph{core-powered mass loss} (CPML) was
introduced by \citet{Ginzburg2016, Ginzburg2018} and refined by
\citet{Gupta2019, Rogers2021}.  \orchard{} implements both competing
forms of the energy-limited cap and selects between them by user flag
\texttt{cpml\_energy\_limit}.

\paragraph{Bondi (sonic-point) rate.}  Following
\citet[][their Eq.~6]{Gupta2019} and \citet[][their Eq.~16]{Rogers2021},
\begin{equation}\label{eq:mdot_bondi}
  \Mdot_{B}
    = 4\pi\,R_{s}^{2}\,\rho_{s}\,c_{s,B}
    = 4\pi\,R_{s}^{2}\,\rho_{\rm rcb}\,c_{s,B}\,
      \exp\!\left(-\,\frac{G\,\Mp}{c_{s,B}^{2}\,\Rrcb}\right),
\end{equation}
where the bolometric (\emph{not} XUV) sound speed is
$c_{s,B} = \sqrt{k_{B}\,\Teq/\mu}$ with $\mu = 2.35\,m_{\rm H}$ for
solar-composition H/He.  The sonic radius is
$R_{s} = G\,\Mp/(2\,c_{s,B}^{2})$.  The exponential factor encodes the
hydrostatic extension from the RCB to the sonic point.  In cool, massive
planets ($G\Mp/c_{s,B}^{2}\Rrcb \gg 1$) it is suppressed by orders
of magnitude, and such planets are stable.  Hot, low-mass planets are
exponentially unstable.

\paragraph{Energy-limited cap (two competing forms).}  An active
debate in the literature concerns which luminosity sets the
energy-limited ceiling.  \orchard{} supports both:

\begin{itemize}
  \item \texttt{cpml\_energy\_limit = gss18}: the form of
        \citet{Ginzburg2018, Gupta2019, Rogers2021}, in which the
        \emph{interior cooling luminosity} provides the energy to
        lift gas out of the well at $\Rp$,
        \begin{equation}\label{eq:mdot_E_gss18}
          \Mdot_{E}^{\rm gss18}
            = \frac{\Lint}{g\,\Rp}\,.
        \end{equation}
  \item \texttt{cpml\_energy\_limit = tang24} (default): the
        re-assessment of \citet{Tang2024}, who argue that the
        radiative atmosphere re-equilibrates on a turnover timescale
        much shorter than the wind timescale.  The \emph{bolometric}
        flux then supplies $PdV$ work in the radiative zone, and the
        cap becomes
        \begin{equation}\label{eq:mdot_E_tang24}
          \Mdot_{E}^{\rm tang24}
            = \frac{L_{\rm bol}}{G\Mp/\Rrcb}\,,
        \end{equation}
        identical in form to Equation~\eqref{eq:mdot_bol}.
\end{itemize}

The composite CPML rate is the smaller of the two channels:
\begin{equation}\label{eq:mdot_cpml}
  \Mdot_{\rm cpml}
    = \min\!\bigl(\Mdot_{B},\,\Mdot_{E}^{\rm chosen}\bigr).
\end{equation}
The two flag settings can yield CPML rates that differ by factors of
3--10$\times$ depending on the planet, and they predict materially
different post-boil-off histories.  \citet{Tang2024} argue their
bolometric form predicts that CPML alone removes only $\sim 0.1\%$ of
envelope mass on Gyr timescales for sub-Neptunes---contradicting the
\citet{Ginzburg2018} picture that CPML produces the radius gap on its
own.  Exposing both forms as a flag lets users reproduce either
literature result.

% ----------------------------------------------------------------------
\subsection{Post-Boil-Off Phase}\label{subsec:postboiloff}
% ----------------------------------------------------------------------

After the boil-off--to--post-boil-off transition (Section~\ref{subsec:transition}),
the planet can contract and cool efficiently.  Three competing channels
are evaluated and the largest sets the wind rate:
\begin{equation}\label{eq:mdot_postbo}
  \Mdot_{\rm post}
    = \max\!\Bigl(\,
        \min\!\bigl(\Mdot_{\rm parker},\,\Mdot_{\rm bol},\,\Mdot_{\rm int}\bigr),
        \;
        \Mdot_{\rm photoev},
        \;
        \Mdot_{\rm cpml}
      \Bigr).
\end{equation}
The ``boil-off-style minimum'' term is included for continuity through
the phase transition; it reduces to the Parker rate alone for
energetically decoupled (radiatively contracted) planets.  In practice
either $\Mdot_{\rm photoev}$ (early post-boil-off, while XUV is
saturated) or $\Mdot_{\rm cpml}$ (late, after XUV decay) dominates.
The $\max$ rule selects whichever channel is currently active.

A global efficiency parameter $\eta_{\rm ml}$ (default 1) multiplies
the final rate in both phases:
\begin{equation}\label{eq:eta}
  \Mdot = \eta_{\rm ml}\,\Mdot_{\rm phase}\,.
\end{equation}

% ----------------------------------------------------------------------
\subsection{Advective Energy Loss}\label{subsec:Ladv}
% ----------------------------------------------------------------------

The escaping wind carries away enthalpy at a rate
\citep[Equation~2 of][]{Tang2024}
\begin{equation}\label{eq:Ladv}
  \Ladv = \frac{\gamma}{\gamma - 1}\,\Mdot\,\cs^{2}\,,
\end{equation}
where $\gamma$ is the adiabatic index (default $\gamma = 7/5$).
This advective luminosity is tracked as a diagnostic but is not
currently fed back into the interior energy budget.

% ======================================================================
\section{Grid Adjustment and Structural Response}\label{sec:grid}
% ======================================================================

After computing $\Mdot$, the Lagrangian mass grid is updated
self-consistently to reflect the loss of envelope material.

% ----------------------------------------------------------------------
\subsection{Uniform Envelope Scaling}\label{subsec:scaling}
% ----------------------------------------------------------------------

The mass lost during one timestep $\Delta t$ is
\begin{equation}\label{eq:dM}
  \Delta M = \Mdot\,\Delta t\,.
\end{equation}
A safety cap limits the fractional loss per step:
\begin{equation}
  \Delta M \le f_{\rm max}\,\Menv\,,
\end{equation}
where $f_{\rm max}$ is the maximum fractional loss per timestep (default
$f_{\rm max} = 0.01$, i.e., 1\%).
A minimum envelope mass floor $\Menv^{\rm min} = f_{\rm floor}\,\Mp$
(default $f_{\rm floor} = 10^{-6}$) prevents complete stripping.

The cell masses $\delta m_{k}$ in the envelope ($k < k_{\rm core}$) are
then rescaled uniformly:
\begin{equation}\label{eq:scale}
  \delta m_{k} \leftarrow \delta m_{k}
    \times \frac{\Menv - \Delta M}{\Menv}\,,
  \qquad k = 0,1,\ldots,k_{\rm core}-1\,.
\end{equation}
The mass-shell boundaries $m_{b}$ are reconstructed by summing the
updated cell masses inward from the surface:
\begin{equation}
  m_{b,k} = m_{b,k+1} + \delta m_{k}\,,
  \qquad k = N{-}1,\,N{-}2,\,\ldots,\,0\,.
\end{equation}

This uniform-scaling approach preserves the relative grid spacing
and prevents mismatches with the hydrostatic and transport solvers
that would arise from removing outer shells.

% ----------------------------------------------------------------------
\subsection{Hydrostatic Re-equilibration}\label{subsec:hse}
% ----------------------------------------------------------------------

With the updated mass grid, the Henyey relaxation solver
recomputes the profiles $P(m)$, $r(m)$, $\rho(m)$, and $T(m)$ by
simultaneously satisfying mass continuity,
\begin{equation}
  \frac{\dd r}{\dd m} = \frac{1}{4\pi\,r^{2}\,\rho}\,,
\end{equation}
and hydrostatic equilibrium (including centrifugal support from
rigid-body rotation at angular velocity $\Omega$),
\begin{equation}
  \frac{\dd P}{\dd m}
  = -\frac{G\,m}{4\pi\,r^{4}}
    + \frac{\Omega^{2}}{6\pi\,r}\,,
\end{equation}
using the Newton--Raphson iteration on $(\ln P,\,\ln r)$ described in the
companion hydrostatic-methods writeup.

The planetary radius $\Rp = r_{b,0}$ (outermost shell boundary)
decreases naturally as $\Menv$ shrinks.  This decrease feeds back into
the next evaluation of $\Mdot_{\rm XUV}$ and the energy-limited rates.

% ----------------------------------------------------------------------
\subsection{Timestep Constraints}\label{subsec:timestep}
% ----------------------------------------------------------------------

When mass loss is active, the evolutionary timestep is additionally
constrained by the mass-loss timescale to ensure numerical stability:
\begin{equation}
  \Delta t \le C_{\rm ml}\,\frac{\Menv}{\Mdot}\,,
\end{equation}
where $C_{\rm ml}$ is a Courant-like safety factor (default $C_{\rm ml}
\sim 0.01$), consistent with the cap in Equation~\eqref{eq:dM}.

% ----------------------------------------------------------------------
\subsection{Small-Timestep Ramp}\label{subsec:ramp}
% ----------------------------------------------------------------------

In runs that begin from a hot, puffy initial state (e.g., $S_{\rm
ini}\!\sim\!9\,k_{B}/{\rm baryon}$), the mass-loss timescale at full
$\Mdot$ can be very short ($\sim 10^{4}$~yr) compared with the
\orchard{} timestepper's initial probe step ($\Delta t \sim
10^{-12}$~Myr).  Applying the full $\Mdot$ at such a small $\Delta t$
produces $\Mdot\,\Delta t \ll \Menv$, but couples a large physical
$\Mdot$ to a numerically tiny step.  Convergence slows, and the rezone
step's hydrostatic re-equilibration is stressed.  Following
\texttt{MESA}'s \texttt{mass\_change\_full\_on\_dt} /
\texttt{\_full\_off\_dt} controls \citep{Paxton2011}, we suppress the
applied $\Mdot$ by a smooth ramp at small $\Delta t$:
\begin{equation}\label{eq:ramp}
  \Mdot_{\rm applied}
    = R(\Delta t)\,\Mdot_{\rm raw},
  \qquad
  R(\Delta t) =
    \begin{cases}
      0
      & \Delta t \le \Delta t_{\rm off}\,,\\[4pt]
      \tfrac{1}{2}\!\bigl[1 - \cos\!\bigl(\pi\,u\bigr)\bigr]
      & \Delta t_{\rm off} < \Delta t < \Delta t_{\rm on}\,,\\[4pt]
      1
      & \Delta t \ge \Delta t_{\rm on}\,,
    \end{cases}
\end{equation}
with
$u = \log(\Delta t / \Delta t_{\rm off}) / \log(\Delta t_{\rm on} /
\Delta t_{\rm off})$.  Defaults are
$\Delta t_{\rm off} = 10^{2}\,{\rm s}$ and
$\Delta t_{\rm on} = 10^{4}\,{\rm s}$.

\paragraph{Self-stabilising rezone near the strip floor.}  The ramp
also stabilises the run as $\Menv$ approaches the floor
$\Menv^{\rm min} = f_{\rm floor}\,\Mp$.  As $\Menv\to\Menv^{\rm min}$
the timestep cap $\Delta t_{\rm ml} = f_{\rm max}\,\Menv/\Mdot$ would
collapse without bound.  In practice $\Delta t$ falls below
$\Delta t_{\rm on}$, $R(\Delta t) \to 0$, and $\Mdot_{\rm applied}\to
0$.  The thermal-evolution cap (\texttt{max\_step}) then takes over and
the run continues without further envelope erosion---this is how the
rock\_gray boundary condition gracefully degenerates to a stripped
super-Earth state when the original sub-Neptune envelope is exhausted.

The dt cap uses the \emph{ramped} $\Mdot$ (not $\Mdot_{\rm raw}$) so
that the cap relaxes consistently with $\Mdot_{\rm applied}\to 0$.

% ======================================================================
\section{Radiative--Convective Boundary}\label{sec:rcb}
% ======================================================================

All mass-loss rate prescriptions require thermodynamic conditions at the
radiative--convective boundary.  In \orchard{}, the RCB is identified as
the outermost envelope cell (cell index $k = 0$).  The function
\texttt{find\_rcb()} returns
\begin{equation}
  \Rrcb = r_{b,0}\,,\quad
  \rho_{\rm rcb} = \rho_{0}\,,\quad
  P_{\rm rcb} = P_{0}\,,
\end{equation}
i.e., the radius, density, and pressure at the top of the model grid.
This identification is appropriate for sub-Neptunes in which the
envelope becomes optically thin near the outermost model cell.

% ======================================================================
\section{Diagnostic Outputs}\label{sec:diagnostics}
% ======================================================================

At each snapshot, the code records the following mass-loss
diagnostics in the HDF5 output (also written to plain-text logs):
\begin{itemize}
  \item $\Mdot$, $\Mdot_{\rm raw}$ --- the applied mass-loss rate
        and the un-ramped, un-globally-scaled physical rate
        [g\,s$^{-1}$].
  \item $\Mdot_{\rm XUV}^{\rm en\text{-}lim}$,
        $\Mdot_{\rm XUV}^{\rm rr\text{-}lim}$ --- the two
        photoevaporation channels of
        Equations~\eqref{eq:mdot_xuv_en} and \eqref{eq:mdot_rr}.
  \item $\Mdot_{B}$, $\Mdot_{E}^{\rm chosen}$,
        $\Mdot_{\rm cpml}$ --- the Bondi rate, energy-limited cap, and
        composite CPML rate of Equations~\eqref{eq:mdot_bondi},
        \eqref{eq:mdot_E_gss18}/\eqref{eq:mdot_E_tang24}, and
        \eqref{eq:mdot_cpml}.
  \item $R_{\rm EUV}$, $K_{\rm tidal}$, $\eta_{\rm XUV}$ --- the CR16
        photoionisation radius, Erkaev tidal factor, and instantaneous
        velocity-dependent efficiency.
  \item $\Menv$, $\Mp$ --- envelope and total planet mass [g].
  \item $\Ladv$ --- advective luminosity [erg\,s$^{-1}$].
  \item $\dot E_{\rm wind} = \Mdot_{\rm wind}
        \bigl[h_{\rm surf} + \tfrac{1}{2}v_{\rm wind}^{2} - GM/R\bigr]$
        --- the wind-energy bookkeeping term of
        Equation~\eqref{eq:Eledger}, used as a global
        energy-conservation diagnostic.
  \item the ramp factor of Section~\ref{subsec:ramp}.
  \item the current evolutionary phase (boil-off or post-boil-off).
\end{itemize}
These permit \emph{ex post} verification of the energy ledger
\eqref{eq:Eledger} and isolation of which physics channel dominates
at any given epoch.

% ======================================================================
\section{Summary of Configuration Parameters}\label{sec:params}
% ======================================================================

Table~\ref{tab:params} lists the user-configurable parameters that
control the mass-loss module.

\begin{deluxetable}{llc}
\tablecaption{Mass-loss configuration parameters.\label{tab:params}}
\tablecolumns{3}
\tablehead{
  \colhead{Parameter} &
  \colhead{Description} &
  \colhead{Default}
}
\startdata
\texttt{mass\_loss}       & Enable/disable mass loss           & False   \\
\texttt{mu\_wind\_amu}    & Mean molecular weight [amu]        & 2.35    \\
\texttt{gamma\_wind}      & Adiabatic index $\gamma$           & 1.4     \\
\texttt{alpha\_kh}        & KH efficiency $\alpha$             & 1.0     \\
\texttt{eta\_mass\_loss}  & Global efficiency $\eta_{\rm ml}$  & 1.0     \\
\texttt{eta\_xuv}         & XUV efficiency $\eta_{\rm XUV}$    & 0.1     \\
\texttt{L\_star\_ml}      & Stellar luminosity [$L_\odot$]     & ---     \\
\texttt{t\_sat\_xuv}      & XUV saturation time [Myr]          & 100     \\
\texttt{min\_envelope\_fraction} & $f_{\rm floor}$              & $10^{-6}$\\
\texttt{max\_dm\_frac}    & $f_{\rm max}$                      & 0.01    \\
\texttt{initial\_phase}   & Starting phase                     & boiloff \\
\enddata
\end{deluxetable}

\bibliography{references}
\bibliographystyle{aasjournalv7}

\end{document}
